\section{从军营到仓场：宋初的中枢为何在开封}

陈桥之后，最先需要解释的不是一个传奇场面，而是禁军怎样回到开封、开封又怎样把一支军队变成一个王朝。后周留下的东京并非空白地面。宫城和官署已有办事的人手，城内外有仓场、桥梁、汴河码头和通往各路的道路，禁军的口粮、军器与俸给须从这些节点穿过。赵匡胤承用这座城，意味着他不必从头制造一套征发和传令机器；同时也意味着新政权被一座需日日供给的巨城约束。《宋史·太祖本纪》与《续资治通鉴长编》建隆元年条能确证即位和朝廷安排，东京城遗址所揭示的城垣、水系和街道使这种选择具有空间感，却不能让我们复原每一所仓屋或每一班挑夫的日常。

禁军的危险正在于它既保护中枢，也可能制造下一个中枢。建隆二年前后，宋廷对石守信等宿将的任官、迁调和优礼，在《宋史》卷二五〇《石守信传》与《长编》的相关条目中留有痕迹。后世把过程压成“杯酒释兵权”，仿佛一席宴饮便消除了军权；较可靠的读法是把它放回数年的任命、调动和饷给重组。将领不再长久以原有部曲驻守一方，军队则越来越由东京的文书、仓储与财政维持。此举降低了五代节度使凭兵力夺位的可能，却把风险移到新的环节：中央若不能准时把粮帛、军器和指令送出，军队的稳定也会受损。

这不是“文官压倒武将”的简单故事。禁军依旧是国家最昂贵也最迫切的资源，文臣、枢密院和转运系统都要与它打交道。南下或北上的一支军队，需要先在仓场取得粮草，再经水陆道路移动；沿途州县既要供应，也要防止军纪破坏接收秩序。国家以轮调限制将领的独立性，往往使远方战役更依赖多个机构的协作。制度解决了谁不该私有军队的问题，却没有自动解决谁能在洪水、歉收或敌军逼近时及时供给军队的问题。此后的统一战争和北方边防，将不断检验这一点。

\subsection{长江、剑门与珠江口：统一不是把颜色涂到地图上}

乾德元年宋军入荆南，高氏政权终结。江陵之所以重要，并不在于它被后世列入“十国”之一，而在于长江中游的城镇、渡口和航道在这里相接。宋军取得江陵，可以使来自东京的命令更容易沿江传递，也能干预荆南原来在后蜀、南唐之间周旋的关系。《宋史》太祖本纪、《长编》乾德元年条与《宋会要辑稿·方域》的改置条，足以说明政权归属的改变；旧官如何留用、运输如何转手、乡村怎样面对新税目，则不会因一条“平定”记录而自动清楚。

伐后蜀使空间的阻力更清晰。964年冬进军，965年孟昶降宋，成都平原与剑门山道同时进入宋廷的账本。平原能供给粮食和人口，山路却限制辎重、信使与伤病者的速度；一座城的降附也不等于整个山区听令。宋方本纪、后蜀世家和《长编》乾德二、三年条提供征服的顺序，但胜利叙事特别容易把军数、俘获和地方反应写得整齐。较稳妥的判断是：宋取得了通往西南的行政接口，随即要面对驻军、旧有地方关系和赋役接续。后来蜀地爆发王小波、李顺起事，不能倒推出965年的每一个人已怀有反抗意图；它却提醒读者，接收从来不仅是换一面旗号。

971年的广州又是另一种接收。珠江口的港埠、盐业、岭南山路和海上往来，使南汉末期的城市与内地州县不在同一运输节奏中。《宋史》南汉世家记录宋军进入广州及刘鋹降服；这能确定政权转换，却不能用中枢的战报概括港口商人、驻军、地方首领或周边居民如何重排关系。宋朝获得的不只是一个城池，也是必须管理的海陆接口。它须将岭南纳入路州体系，又不能假定来自东京的规章能够毫无折损地抵达每一处港湾和山谷。

974年渡江围金陵、975年李煜降宋，则把江防、船只、城中储粮和江南财赋聚到一场长期的压力里。南唐都城并非静待接收的容器；围城中的军民、沿江运输者和周边州县都被战事改变。宋方《宋史·南唐世家》与《长编》开宝七、八年条为攻围和降服给出年序，却不足以量出伤亡、饥馑或每处财赋恢复的速度。宋廷此后能调动江南的资源，正因它逐步把水网、仓储和官员接入自己的制度；但“逐步”不能被删去。统一使中央拥有更宽的供给面，也使它必须对更广阔、差异更大的地区承担行政成本。

\subsection{河东之后的北方：没有缓冲政权的道路}

976年太祖去世，太宗即位，继承的是南方接收尚在展开、河东和燕云仍在宋境之外的局面。关于继承的后起传闻不能补作宫门内的对话；《宋史》两朝本纪和《长编》开宝九年条所能确认的是权力已经转换。979年五月宋灭北汉，太原进入宋朝控制范围。北汉长期与辽相连，因而太原的改变也改变了宋辽之间的距离：原先的缓冲政治体消失，宋军和辽军的道路、堡寨与侦察线更直接地相遇。

太宗随后北攻幽州，在高梁河受挫南撤。《宋史》太宗本纪、《辽史·景宗本纪》与《长编》太平兴国四年六月条都记录了这一节点，却对兵力、追击、皇帝处境和战功各有叙述。将它们并列，不是为了挑出一个绝对中立的数字，而是为了看见双方各自的战争场景：宋军希望在河东既定后夺取燕云，辽军则能依托近畿、骑兵机动和既有防御网络应对。败退的直接后果是幽州未得；更长的后果是东京必须为河北、河东的防务、粮道和驻军做持续预算。

986年雍熙北伐以多路军北进，仍未改变燕云格局。三路协同行动的设想本身说明宋廷希望用更大的文书和补给体系补偿边境距离；它也放大了消息延迟、行军季节、将领协同和粮道安全的风险。宋、辽史书对胜败和责任的书写都带有战后政治的颜色，因而不宜把某个军数或单一将领的动机写成定论。可以看到的是，此后“守边”不再只是几座关城的问题：河北农户、沿线州县、转运官和驻军都被卷入一个需要经常供给的北方体系。

宋初限制将领独立性，在这种远方战争中显出两面。它防止地方统帅凭一军成为新的割据者，却要求遥远的前线等待更复杂的中枢协商。粮食和军器经由州县、河流和道路转送，任何一个节点的拥塞都可能改变军队的行动。不能把高梁河和雍熙的失利机械地归因于“重文轻武”；该说法遮蔽了燕云地形、辽的兵力组织、宋军补给与指挥网络之间更具体的限制。制度的得失在这里很明白：它使中枢免于五代的兵变恐惧，却使北方防线成为一种长期、昂贵而脆弱的联结。

\subsection{澶州的寒冬：盟约如何进入账簿和边市}

997年真宗即位时，宋辽之间并未因此前战事而停止来往。使节、边市、侦察与堡寨构成一种既对抗又交涉的边境日常。1004年辽圣宗与萧太后南下，宋军在澶州附近应对，真宗至前线；《宋史》真宗本纪、《辽史·圣宗本纪》和《长编》景德元年冬条提供相互牵制的年序。它们无法使我们听见每次军议，也不宜据某一方的胜利语言判定谁完全控制战场。澶州之所以成为转折，是双方都要把继续作战的代价放在眼前：渡河、城防、骑兵、粮食、冬季与国内政治同时限制选择。

1005年达成的盟约以银绢、使节和一套称谓安排维持停战。宋方材料常把支出写作“岁币”，后世又常把这一词化作民族胜负的判断；《宋会要辑稿·蕃夷》的条目与《辽史》相对照，更能说明它是一种把战争风险转换为可预期财政支出的制度。银绢需要从财政网络中筹出，使节需要沿道路安全往返，边市和堡寨仍须维持。盟约既不是辽对宋的臣属，也不是宋买来的一劳永逸和平；它是两国在军事与资源限制下建立的程序。

盟约签订后，边界也没有变成一条无人地带。1010年前后辽军南入和宋廷的河北防务安排显示，纸面关系要靠地方的侦察、修寨、粮运和信息交换才可能持续。对边民而言，和平常常意味着另一种日常：军役未必消失，市场和通行可能受到新的管束，地方官还要在警戒与交易之间判断风险。现存朝廷材料比边民自述多得多，故不能把盟约描述为所有人共同感到的安定；但它确实把北方竞争从周期性的大战转入有使节和财政程序可追踪的长期关系。

真宗后来举行封禅，把边境暂稳译作皇权与天命的图像。《宋史》与《宋会要辑稿·礼》可让我们看到仪式的筹备和举行，却不能把祥瑞记述当成国家强盛的独立证明。对东京而言，礼仪展示的是被组织出来的秩序；对河北而言，秩序仍要靠仓储、兵员和道路维持。这两种尺度并不矛盾，只是不能混为一谈。

\section{粮食、纸券与城市：国家从哪里取得可用的资源}

北宋的财政不能被一句“富庶”概括。开封的禁军、河北的防务、河工、官员俸给和使节往来，都要求把不同地区的实物、钱币和劳力接到同一网络中。网络越大，地方差异越重要：江南水网的粮食、四川盆地的市场、盐场和冶场的收入、边区的军需，不能用同一条道路或同一类文书处理。制度材料能告诉我们某项政策被怎样命名、何时颁行；它不能单独证明一户农家实际交了多少粮，或一个渡口在雨季仍能通行。

\subsection{成都的起事和纸券：同一盆地的两种压力}

993年王小波起事，994年李顺一度占据成都，宋军随后平定。《宋史》卷四六四《王小波传》、太宗本纪与《长编》淳化四、五年条保留了起事的中枢视角，其中“均贫富”的口号尤其醒目。它不能被当作参与者社会构成的完整调查：镇压者的叙事会分类首领、兵众和罪名，地方普通人的选择却较少留下文字。仍然可以看出，蜀地在新朝接收后并不是安静的财赋仓库；土地、税役、军政和区域权力的压力能够在成都平原集中爆发。把事件只当作“乱民”插曲，便会抹掉统一国家如何在地方被感受和抵抗的问题。

数十年后，益州交子务的设置又显示四川的市场和货币问题。1023年官府接管交子，相关材料见《宋史·食货志》及《宋会要辑稿·食货》；它不是纸币忽然取代铜钱的神话起点。商人原有的信用票据、铜钱供应、兑现能力和地方贸易条件，才是官府介入的背景。设务使官方能够规定发行、收换和管理的接口，也把信用风险带入行政文书。发行额、能够兑现的数量和实际流通范围须分开判断；不能由一条制度记录推及全国货币化，更不能将后来的纸币问题提前写入1023年的成都。

起事与交子并置，并非要把它们归入一个抽象的“商业社会”。前者使地方接收、赋役和暴力的紧张显形，后者让我们看见商人、钱铺与官府如何在一座区域城市处理交换的不便。它们都提醒读者：财政不是从中央向外单向流动的数字，而是许多人在田地、市场、仓场与衙门之间不断协商的关系。档案保存偏向官府与有书写能力者，故我们仍不能替成都所有居民说出同一种经验。

\subsection{稻种、田地和水网：增长不是一条诏令}

1012年前后朝廷自福建调入占城稻种，史料常被用来说明江南农业的变化。《宋史·食货志》、真宗朝编年材料和后来的农书可追索这一过程，但它们各自记录的是朝廷的关切、技术知识或后设总结。稻种能否在某处成功，取决于水利、土壤、劳力、租佃关系和农户能否承受试种风险；一份调种记录不能直接换算为全国增产率。更何况江南也不是同质空间：湖泊、河港、堤岸和村落的距离，会改变一场洪水或一段旱期如何落到家庭身上。

粮食一旦被征入国家体系，又要面对储藏和运输。汴河把东南与东京接起，但水道的便利不是自然赠予：闸、堤、河工、船户、仓官和沿线州县都要参与，淤塞和水位变化会使看似稳定的供给变得紧张。河工和漕运的制度条目告诉我们国家不断修治和调度，却不能让我们断言每一项工程都达到预期。对开封的禁军而言，水路上的延误是饷给问题；对沿河居民而言，它也可能意味着征夫、占地、堤防风险和临时劳役。把这些后果分开，才能避免把“漕运发达”写成无人的基础设施。

农村家庭的材料尤其零碎。墓志、契约和地方文书更多保存有产者、能够立字者与被卷入诉讼者的经历；妇女、佃户、短工和迁徙者往往只在亲属或官府分类中短暂出现。因此，本书不以一个“宋代农民”的形象代替不同地区的家庭。可以较有把握地说，国家的粮税、劳役和赈济会进入地方生活；至于每次征收如何在某村分摊，须等待具体文书和区域研究的支持。

\subsection{汴京的门、桥和夜市：城市不是单纯的消费舞台}

开封的繁华常被《东京梦华录》的街市、酒楼和节令描写所代表。孟元老写作于北宋覆亡后，怀旧的城市笔记使许多店铺、瓦舍和道路变得可见，也使读者容易忘记它不是一份逐日城市统计。它最能照亮的是某种记忆中的都城景观；城址考古、钱币、制度条目和其他文本须共同限制它的范围。即使如此，汴京的城市生活仍不能只从皇宫看：门禁、桥梁、河道、货栈、军营和市场把官员、商人、工匠、僧道、脚夫、妇女和外来者置于相互看得见的空间。

城市市场所需的不只是货物。度量、税卡、巡检、消防、治安和诉讼都让交易与行政相遇。白天的市与夜间的街并非简单的自由／管制对立；某些地方允许的买卖，可能在另一类货物、另一段时间或另一处城门受到约束。官府对市税和价格的记录反映其征管兴趣，笔记文学则偏爱奇观和消费，两者都不能自行代表全城。工匠、雇工和小贩常较少有可存的自述，不能因材料沉默就把他们从城市的生产与服务中删去。

东京也是巨大的财政装置。禁军和官员的消费能够支撑市场，市场和税收又反过来帮助供养中枢；一旦战事、河运或货币信用发生波动，这种互相依赖也会传递压力。城市并非“农业剩余自然汇集”的结果，而是道路、水道、制度和劳力持续维持的关系。北宋后期的花石运输和靖康围城，将从不同方向暴露这种关系的脆弱性。

\section{边寨、账簿与地方首领：仁宗朝的不同边界}

北宋的边防不只在河北。西北面对的是党项集团逐步建国，西南面对的则是山地、水道和地方首领之间难以用州县名目概括的关系。仁宗朝的朝廷争论常把它们归为“边事”，但泾原路的堡寨、广西的洞寨、开封的财政和地方社会承受的压力并不相同。将其并列，是为了让读者看见同一个中枢如何在不同地理中遇到不同的行动者，而不是把所有边境写成等待命令抵达的边缘。

\subsection{元昊称帝：宋廷并未拥有命名边界的权力}

1038年李元昊称帝建国，西夏由此以自己的政治语言进入宋辽之间的关系。《宋史·夏国传》和《续资治通鉴长编》记录宋廷的反应，称谓中带有宋朝的正统判断；要理解事件，不能让“僭号”一词替代党项国家形成的过程。河西、灵州、兴庆一带的城镇、绿洲、牧地和道路，使西夏能够在宋辽之间调度人力、贡赐与军队。宋方材料对于西夏内部的财政、部族关系和地方经验并不充分，故本章只在可证范围内说明双方接触的军事与外交后果，不把宋廷的命名当作对方社会的自我描述。

宋廷最先面对的是现实的军事成本。西北堡寨需要粮、草、军器和轮换士兵；一条从京师发出的命令，到达环庆、泾原或鄜延时，已经经过转运、仓储、道路和许多地方人员。战报擅长记录城池、斩获与将领任免，却常夸大己方战果、压缩普通军民的处境。兵额、伤亡和骑兵数若无双方材料或地方遗存限制，不能写成精确事实。可以肯定的是，西夏建国使宋的北方财政不再只围绕辽，而是要在更广的边线配置资源。

1040年代宋夏战争的城寨攻守，说明战略不是地图上的箭头。谁能守住渡口、山口和水源，谁就可能决定粮车能否通过；雨雪、牧草与道路季节性也会改变行动。宋军未能以一次战役消除西夏，西夏亦不能把宋的资源网络排除在河西以东。双方逐渐在战争、互市、册封和赐予之间转换，这种转换并非和平的反面，而是边境社会继续生活的条件。边市若能开通，商人和牧民才有交易机会；边市若被关闭，州县、军寨和周边人群便要承受不同的价格与供给压力。材料很少直接记录他们的声音，不能把朝廷协议写作人人共享的结果。

\subsection{庆历更张：朝廷用学校与考课回应边防压力}

1043年范仲淹、富弼等所推动的庆历新政，常被后人放入“改革失败”的短句里。就当时而言，朝廷讨论的是官员考课、学校、选任和财政等问题，背景正包括西北军费和中枢对行政效率的焦虑。《宋史》相关列传、《长编》庆历年条和《宋会要辑稿》的职官、学校条目，能让我们追索某些诏令与机构安排；它们不能证明每一所州县学校立即办成，也不能证明一个考课制度已改变全部地方官的行为。

庆历改革的困难不是某个人是否“有道理”，而是国家如何把政策变成地方可执行的动作。学校需要经费、师资、场所和学生来源；考课需要能够上报并核验的文书；边防所需的粮饷又与这些项目争夺财政和官员时间。制度如果反复申令，往往说明执行存在摩擦，而不只是表明中央勤政。庆历后党争和人事变动削弱了改革的连续性，但“失败”并不意味着它毫无遗产：学校、选任和对地方行政的关切仍会在后来的争论中以不同形式出现。

\subsection{贝州与河北河患：边防地带也是日常生活的地带}

1048年贝州军卒起事，王则据城，宋廷调兵平定；同年黄河商胡决口，河北的堤防、田地和道路又受到冲击。《宋史》仁宗本纪、王则传与《长编》庆历八年条使这两组事件进入同一年度，却不应匆忙把它们编成一个唯一的因果链。军队不满、地方控制、河道变化和灾后供给各有条件，材料也各有偏向。将它们放在一起，是为了看见边防州县并非只有“前线”这一面：那里有军营、户籍、农田、堤防、市场和家族，任何军事或河工命令都可能穿过这些关系。

黄河决口的报告常以河道、受灾州县、工程与赈济的行政语言出现。它能够说明官府在何处看见险情、如何调度工程，却不能仅凭一条奏报计算死亡人数或复原每一户的迁徙。堤防工程需要征夫、材料和地方协调；在边防区，河工又可能同军粮、修寨争夺劳力。河北居民面对的不是抽象的“国家能力”，而是何时能得到粮食、哪条道路仍可走、田地是否被水淹没、役夫是否必须离家。存世材料多由官员书写，故这些问题只能以有限证据提出，而不能虚构灾民的共同心声。

\subsection{侬智高与广西水道：州县之外的政治关系}

1052年侬智高起事，一度攻入邕州并向北推进，后由宋军平定。宋方本纪、列传和《长编》记下军事行动，却把复杂的边地关系压缩为叛乱与平乱。广西、交趾边境的山地、河谷和洞寨并不以开封的州县格局为唯一秩序；地方首领、贸易路线、婚姻和兵力动员形成的网络，决定了谁能利用水道、谁能在山间补给。不能把侬氏及其追随者简单写成一团固定的“族群”，也不能让宋朝官修的罪名替代其行动条件。

狄青南下的军队要通过岭南道路和水道，军事胜利取决于向导、补给、季节和地方关系，并不只是朝廷授予一纸命令。战后宋廷加强控制的措施，也不能自动证明边地从此完全纳入中枢。州县设置、军寨和使者往来能够改变资源流向，但洞寨社会和跨境联系不会因此消失。这个事件提示读者，北宋的“边疆”不是一条从东京向外递减的线：西北与辽夏相接，西南则与交趾及地方网络相连，国家须以不同方式处理道路、军事和财政。

\subsection{从边事到财政：一项解决办法留下什么压力}

仁宗朝以和议、堡寨、学校、任官和军费应对多重压力。它们各自解决了部分问题：宋夏停战降低了持续大战的风险，边寨组织了守备，庆历政策试图改善行政接口，南方平定恢复了部分州县秩序。但每项安排都有成本。银绢和军饷来自财政网络；堡寨需要长期补给；考课和学校需要地方人手；战后控制要求重新分类人群和道路。没有哪一项可被称为单纯的成功或失败。

这正是北宋后来改革会如此尖锐的原因。熙宁朝的官员并非面对一个未曾治理的国家，而是面对一套已能征发、转运、交涉却仍被边防、河工和地方差异牵动的制度。新法将尝试改变财政和基层接口；它的执行也会把更多书手、州县官、富户、借贷者和服役者卷入争议。

\section{熙宁的试验：把财政能力送到州县，还是把州县送进账簿}

1067年神宗即位，次年王安石入参政，1069年青苗法颁行。时间节点可由《宋史·神宗本纪》《长编》熙宁年条和《宋会要辑稿·食货》追索；政策被写成一串名称时，最容易失去其实际问题。中央所关心的是军费、财用、流通、劳役和边防的承受力，改革者希望减少中间耗损、使资源更可计算；反对者担心地方官的考成、强制和市场干预会把原本可协商的关系变成新的压力。双方所见的不是同一个州县，也不能把他们的言辞直接等同于每个借贷者或役户的经验。

\subsection{青苗：春天的种子、秋天的催纳}

青苗法规定由官府在青黄不接时向农户提供钱或谷，秋后归还并计息。它试图在民间高利借贷和国家财政之间建立接口：若基层能够以较低成本取得周转，国家也能使闲置资金流动并获得收入。制度文本能够说明设计的方向，却不能证明每一笔贷款是否自愿、利率如何被理解、歉收时是否宽缓。不同州县的土地、作物、富户结构和官员考核并不相同，同一法令在不同地方会遇到不同后果。

争议的核心不应被写成“救民”与“害民”两句口号的对决。地方官要编定户等、登记、发放和催纳；富户可能担任保人或影响分配；贫困家庭则可能在短缺、税期和劳役之间作选择。官府文书常保存应当如何办和某些官员的奏报，反对派文集则强调弊端，二者都带有立场。若无当地契约、诉讼或其他材料，不能替某一村庄判定全体人的感受。能够确定的是，青苗让财政管理深入触及季节性借贷，并使国家责任与地方压力的界线更难划分。

\subsection{免役、市易与均输：市场成为行政问题}

免役法以钱代替一部分差役，再由官府雇募人员服役，意在把不均的差役负担转为可计算的财政项目。可是“以钱代役”不意味着负担自然公平：谁被定为出钱者，谁能得到雇役机会，地方如何估价，都需要具体的户等、文书与执行人。制度若只看中央条文，容易显得整齐；从州县看，它要求更多登记、催纳与申诉，也使原有的人情和地方权力以新的方式进入行政。

均输法和市易法则把运输与交易直接纳入政府经营。均输意在利用各地价格差和转运渠道降低采购成本；市易务以官本参与货物周转、抑制大商人的操纵。开封的市场、仓库、船运和度量衡使这类设想有现实依据，但价格并非由一处衙门单独决定。货物的季节、道路、商人信用、地方税费和信息速度都会改变实际结果。市易文书和反对者奏议都不能自动代替市场参与者的全部声音；更不能由“官本”二字推论国家已经控制所有交换。

改革的代价亦在于文书。每项法令需要能够识字、计算、登记和追催的人；考成制度若以完成额数为重，地方官就可能把复杂的生计关系压成可上报的数字。反对者批评的“扰民”往往指向这种压力，改革者则认为旧制度已让中间人侵蚀公私利益。两种观察都有材料基础，却都不能免除区域差异。北宋的改革并非抽象的国家与社会对立，而是国家通过更多地方接口进入社会，社会也借这些接口提出、规避或重写自己的要求。

\section{从熙河到元祐：边功、法令和朝廷的时间差}

新法并非只在开封的堂议中运行。1060年代末至1080年代，神宗朝将财政改革同西北扩张相连，熙河路的开辟便是最明显的例子。王韶在洮河、河湟一带与吐蕃诸部、地方城寨和宋军发生的接触，不能被概括为“收复失地”；这里既有不同语言和宗教传统的人群，也有山谷、牧地、道路和粮运的限制。《宋史》神宗本纪、王韶传、《长编》及会要的边防条可勾画宋廷任命、筑城和置路的年序，不能据此替各部说明政治意图，更不能把宋方战报的户口、俘获和地界写成精确、稳定的控制范围。

\subsection{熙河路：一座城寨要靠什么留在地图上}

熙河经营的现实问题是补给。军队进入河湟，需有能够通行的山路、可驻扎的节点、牲畜和粮食；新置的城寨若没有居民、驻军、市场和转运，便只是文书上的地名。宋廷希望以城寨、互市和招抚切入高原与河谷之间的网络，地方行动者也会衡量宋军、吐蕃部族与既有贸易关系带来的利害。官方材料常将招纳写作某一方归附，实际关系却可能随季节、军力和资源变化。因而本章只说宋的行政和军事影响有所伸入，而不以一个“归附”词抹平当地的多重选择。

边功与财政相连。筑城需要木石、劳役与军器，驻军需要粮草，远距离运输又会损耗牲畜与人力。神宗朝的改革者认为整顿财用可以支持边防；批评者则担心开边会使原已紧张的州县和仓储负担更重。两边都不能仅以道德名词理解：某一座堡寨的战术作用，可能给相邻地区带来征发；一项看似成功的招抚，也可能因后续供应不足而难以维持。北宋的“富国强兵”因此不是先富后强的直线，而是军事、行政和地方承受力彼此拉扯的过程。

\subsection{元丰改制：把职责写清，不等于把事情做完}

1080年前后，元丰改制重整中枢官制，试图以三省、六部等名目厘清职责。制度史常将它写成对唐制的复归或对宋初体制的校正；对当时官员而言，改制还意味着文书须重新流转、职位须重新安置、部门间的责任须被定义。《宋史·职官志》《宋会要辑稿·职官》与《长编》能记录诏令和名目，却不能凭一纸制度设计证明所有机关立即顺畅协作。新的名称也不必然等于旧的职能，尤其在财政、军政和地方通信相交的地方。

元丰年间的官制整理与新法同样面对执行问题。中枢可以规定奏报由何处经过，州县仍要有能处理文书的人；中央可以规定职掌，边防和河工的突发事务仍可能越过清晰的分工。改制的优点是把权责和程序做得更可见，代价则是转变期间的磨合与职位竞争。后世若以“恢复祖制”或“破坏祖制”概括，便看不见制度如何在具体文书、人员和地方需求中被使用。

1085年神宗去世，哲宗即位，太皇太后高氏临朝，司马光等人主导废止或调整部分新法。此后的元祐政治不是一支“旧党”自然取代一支“新党”：青苗、免役、市易等法牵连的机构、收入和地方关系已经存在，废改本身也要依靠诏令、官员和州县执行。新旧政策的来回改变了人事与言论空间，却不能让此前的借贷、差役、市场和军费问题消失。将派别称谓直接投射到所有地方官和民众身上，只会把复杂的行政变动缩成首都的道德剧。

\subsection{河流、灾荒和赈济：国家看见灾害的方式}

北宋的河工与赈济材料十分丰富，因而也容易制造一种国家无所不见的错觉。官员上报决口、旱涝、虫灾或饥民，是因为这些现象进入了行政视野；报表中的户数和粮数又服务于请粮、减税、修堤或考课。它们首先说明某机关在某时以某种单位看见了问题，并不自然等于实际受灾人口。黄河改道、堤防工程和淮河水网的变化还需要地理与考古研究限定，不能用一条纪事建立全境环境史。

赈济也是一种道路问题。粮食若在仓中而道路中断，灾区仍不能得到；若地方官没有人手登记，减税和赈贷也可能无法按条文分配。灾民并非被动的数字，他们会投靠亲属、迁往城镇、变卖财物、改种或离开原籍，但这类策略只有在个别契约、墓志、笔记或地方记录中偶尔可见。我们不能替他们编造一致的遭遇。可以说的是，河流和气候迫使国家的征税、兵粮和城市供给不断重新安排，北宋财政的稳定从未与环境风险分离。

\section{文字、学校和技术：谁能进入国家的视野}

北宋的政治争论之所以密集，与识字者、科举、印刷和文集形成的沟通网络有关。但“士大夫社会”不等于全体社会。榜单、奏议和文集最清楚地记录了能够应试、任官或写作的人；工匠、雇工、妇女、僧人和乡村多数人的生活，往往只通过契约、法律、墓志或他人的叙述片段出现。文化史若只追随名臣的文章，就会把支撑文本生产的纸墨、刻工、书肆、寺院和家庭劳动变得不可见。

\subsection{科举的门槛和书院的地方性}

庆历以来的学校政策和嘉祐科举，使教育成为朝廷讨论治理能力的一部分。学校条目、考试制度和登科录能说明国家怎样定义合格的知识；它们不能说明每个州县都有同样师资，也不能把一张榜单转换为“社会开放”的总指数。准备考试需要时间、书籍、教师和家庭支持，因而入仕渠道虽比门阀世袭更具流动性，仍不可能对所有人平等。地方家族投资教育、寺院和书肆保存或贩卖文本，都是把抽象的考试制度变成实际机会的条件。

苏轼、司马光、王安石等人的文章与政论，留下北宋政治最清晰的自我表达之一。它们能让我们看见官员如何描述财政、刑罚、学校与边防，也会以修辞塑造对手。将其当作辩论的一方证词，而非对社会的透明记录，才能避免让名臣的语言覆盖其他行动者。乌台诗案中苏轼受审，正显示文本、监察和派别如何在国家机器中相遇；案件材料能说明政治压力，不能用来推定所有读者如何理解一首诗。

\subsection{仪象、历法和城市工匠}

1092年苏颂主持的新仪象台完成。技术史常以复杂的机械装置展示北宋知识成就，但一座仪器的意义不在于替后世发明史提供孤立的“领先”证明。历法、授时、天象观测和朝廷礼仪相连，官府需要工匠、材料、维护和能够解释数据的官员；仪器设计被记录下来，也不等于实物此后持续完好运行。《宋史·天文志》、苏颂的技术著作及后来的图像传统可以照亮不同层面，却不能把传世图纸当作完整现场。

工匠在此并非无名的背景。铸造、木作、纺织、造船、印刷和陶瓷生产都要求熟练劳动、原料来源、市场和管理。窑址或器物能显示某处生产和流通的痕迹，不能替整个行业估计产量或劳动条件；官府条文能规定役使和采购，也不能自动说明每一次执行。把技术放回城市、财政和劳动关系，才会看到仪象台和书籍一样，依靠许多并未留下个人文字的人。

\subsection{寺院、功德和公共空间}

寺院是北宋城市和乡村中保存文字、接待旅行者、组织仪式和拥有财产的场所，但不能被写成一个脱离国家与社会的统一机构。僧道名籍、敕令和寺院碑刻各自呈现不同事实：官府关心度牒、管理和财政，捐施者关心功德与家族记忆，碑刻又常用理想化的语言展示秩序。寺产、香火、教育和救济在不同地方的作用需要分别证实，不能由一通碑文推及全国。

城市中的寺观还与市场、道路和节庆交错。它们可能提供识字、借住、施舍或交易相遇的空间，也可能成为征发、禁令和政治监督的对象。如此理解宗教，不是把它降为财政工具，而是承认信仰、制度和物质资源在具体地点相互影响。对于无名信众，我们通常只能见到物件、题记和后人的分类；材料的沉默应保留在叙事中，而非用想象的仪式和内心填满。

\section{宣和的转运：宫廷审美如何遇到江南道路}

徽宗朝的宫廷画院、礼制与道教活动，使北宋后期的文化形象极为鲜明。然而政治与财政压力并没有因为艺术活动而停止。1100年徽宗即位后，蔡京等人在财政、官僚与仪礼上发挥影响，朝廷又在东南搜求花石。相关本纪、会要和笔记保存大量名称和批评，材料层次却不同：官修史倾向以亡国后的政治评价组织人物，笔记易将奇闻和暴政结合。可以确认的是，运输与征发在两浙等地确实成为问题；不能把后来所有民变都简单归因于一支花石纲。

花石从太湖和江南地区向东京移动，需要船只、河道、役夫、地方官和沿线停靠点。对宫廷而言，它们是园林和权威的物件；对承担运输的人而言，它们可能意味着占用舟船、道路和劳力。正因两种视角不一样，花石纲是观察北宋末年资源配置的好入口：中枢能把远方稀有物集中到城市，说明它仍有调度能力；这种能力若无视地方生活和既有运输需求，又会制造新的冲突。技术上的可行不等于政治上的可承受。

1117至1121年方腊起事发生在这一更宽的背景中。宋方史书和后来的叙事记录了起事、攻守及镇压，常以宗教、盗贼或妖术为分类；这些称谓能反映官府如何处理威胁，却不足以说明参与者共同的信念和社会构成。东南的税役、劳力、市场与地方权力各有作用，不能用一条“花石纲逼反”取代。方腊事件使国家必须将军队和资源从其他方向调回，也显示两浙水网与山地并不只是供应东京的后方，而是会以自身的矛盾改变全局。

\section{海路、燕云与金：新的联盟怎样改变了旧的边境}

女真诸部在1115年建立金朝，辽的东北与北方秩序随之变化。对北宋而言，这不是一条来自远方的新闻，而是重新触及燕云问题的机会与风险。宋廷长期以澶渊关系处理辽，燕云仍在辽的控制下；金的兴起使宋朝能够想象以新的同盟取得旧日未得的城地。这个判断本身已经把外交、海运、军队和燕地居民放在一个选择中。《宋史》徽宗本纪、金国传和《三朝北盟会编》保留宋方的交涉线索，《金史》则显示金朝自身的行动年序；双方材料都服务各自政治叙述，不能据此推知每个使者或边民的内心。

\subsection{从登州出海：联盟不是在地图线上签署的}

1118年前后宋金使节开始接触。使节取道登州、横渡渤海再向女真方面行进，这一行程使联盟首先成为海路、船只、港口和翻译的问题。渡海并非把两国之间的辽朝空间简单跳过：出发者需要获得船只、等待风向、携带文书和礼物，还要让港口与沿岸人员配合。宋方编年材料记载遣使和答使，能确证交往的时间；它很少保存船员、地方商人和译者的完整视角。把海路写成一条清晰的外交箭头，会抹去使团实际依赖的城市和劳动。

宋朝希望金军攻辽时取得燕云，金朝则以自己的军事和资源计算同宋交往。1120年海上之盟的条款、称谓与互相承诺，常被后世读成一方的高明或失策；在当时，它是两个力量不同、信息并不对称的政权试图安排共同敌人的文书。盟约能够规定目标，不能保证军队到达、地方接受或资源足够。辽在燕云经营已久，城镇官民、道路和粮仓不因宋金协定而失去作用；任何接收都要面对他们的选择和战争带来的破坏。

联盟给北宋带来一个难题：它一方面可借金牵制辽，另一方面又必须让金了解宋在北方的军力、财政和战略意图。使节来往和礼物馈赠本是外交程序，也会成为信息流动的渠道。不能把这一点写成宋方“泄密”或金方“欺骗”的单一心理故事；各方在不确定中使用有限的信息作判断，后来结果并不使当时所有选择天然荒谬。燕云的价值正在于它不是抽象的旧疆，而是与河北平原、关隘、城镇和通向开封的道路相连的空间。

\subsection{攻辽与燕京：取得城池不等于取得防线}

1122年宋军北伐辽，在燕京一带受挫，金军随后取燕。宋方《宋史》和《北盟会编》、金方《金史》均记载相关行动，但对指挥、兵力和功劳的分配不同。宋军的问题并非一个将领勇怯能够概括。军队远离稳定补给线，须穿过河北的道路并面对辽的守备；金军则已在东北和燕云的战事中拥有自己的机动和目标。战役显示宋廷仍无法仅凭中枢命令迅速把北方资源组织成可取燕云的军力。

燕京转归金方后，宋廷以财物和交涉取得部分燕云州县。条约上的“交割”看似明确，实际却意味着城门、仓库、户籍、军队、官员和居民关系要被重新分类。谁愿留、谁被迁、谁能继续经营市场或耕种土地，都不是盟书本身能规定的。宋方得到若干城地，却也得到一条更难保障的前沿：原来隔着辽的宋金，现在在更靠近河北腹地的地方直接相对。将这段写成“收复燕云”或“得地失地”的二分，无法说明防线为何反而更不稳。

辽亡并未使北方自动和平。1125年金灭辽后，原有的使节、边市与军事均衡被打破。对金朝而言，宋是否履行联盟中的协助、能否守住新得城地，成为判断其行动的因素；对宋廷而言，金军迅速扩展也使原先的伙伴变成潜在威胁。责任归属在宋金两方史书中各有叙述，现代研究也仍讨论盟约、交涉和军事部署的相对作用。本书所能确定的是：宋辽之间经百年程序维持的边境，已被一次新的联盟重置，而新的安排没有建立起等量的互信、补给与防御能力。

\subsection{第一次围城：河北道路如何把战争推向东京}

1125年冬金军南下，宋廷在河北的防线迅速受压。军队能否集结，取决于道路、粮食、城寨、河渡和地方官的判断；勤王军的名单和战报在史书中十分醒目，实际到达的速度、人数和粮饷却难以由一方材料完全核对。宋朝长期依赖东京的转运和命令网络，此时必须在敌军逼近、地方恐慌与朝廷争论中同时运作。中枢的诏书能传出，不等于每一处州县都有能力执行。

1126年金军第一次围攻开封，宋廷以议和与财物使围城暂时解除。围城使都城的双重性质完全显露：它既是政治象征，又是人口、粮食、军器和信息高度集中的城市。城门一旦受威胁，官署的决策、汴河的供给和市场的价格都会相互影响。《宋史·钦宗本纪》《北盟会编》及金方材料能确认围城和议和的轮廓；关于城内恐慌、某人言辞或具体财物数额，叙述间多有差异，不应被戏剧化为统一的城市体验。

暂时解围并没有恢复旧秩序。议和所需的财物要从财政和城市资源中筹出，军队的缺额、河北的失守和各地消息的迟滞仍在。更重要的是，金宋关系已经从有共同目标的联盟转为直接的军事与交涉关系；此前与辽相处时形成的澶渊程序不能直接移用。北宋在第一次围城后仍可能作出不同的军事、外交和财政选择，但每种选择都受到时间、资源和信息的限制。把后来城破视作此刻的唯一结局，会抹去当时仍存在的犹疑和行动。

\section{靖康：一座城市失守，许多秩序被迫分开}

1126年末至1127年初，金军再次南下，开封陷于更严峻的围困，最终城破。靖康是北宋叙事的终点，却不应把它写成一个只属于宫门的瞬间。《宋史·钦宗本纪》《北盟会编》记录朝廷、使节、守城和迁徙的线索，金方《金史》也保存征服一侧的年序；它们都集中于皇帝、将领和议和，城市中普通居民、工匠、商贩、妇女和流动人口的经历则难有连续档案。承认这种不均衡，比为所有人编造同一种灾难体验更接近材料。

\subsection{城内的粮食、钱币和信息}

围城首先改变的是流通。城外道路受阻，粮食、燃料、军器和消息都难进入；城内存粮、仓库、市场和军队随即成为焦点。价格变化与饥饿在后来的叙事中经常相连，但具体物价、存粮和死亡数字须看记录来源和场景，不能用一个令人震惊的数字替代城市经济的全部状态。官府会优先考虑守军、宫廷和秩序，居民则要在亲属网络、市场和有限的救济之间寻找办法；谁有储藏、谁能接近官府、谁被困在城外或城门附近，处境绝不相同。

信息同样是一种资源。朝廷收到的战报、使节带来的条件、城外传闻和街市流言不会完全一致。官员围绕守战、和议和任将的争论，不只是性格冲突，也来自对敌军位置、援军可能和财政能力的不同判断。后世史书以结局组织这些意见，很容易让某些主张显得天然正确；在当时，作决定的人却不能知道后来发生的一切。我们可以评估制度和战略造成的约束，却不能把未知的结果倒灌为人人早已知道的命运。

\subsection{被带走的人与留下的地点}

1127年金军攻入开封，徽、钦二帝及宗室、官员等被北迁，北宋中枢至此终结。被俘者名单和朝廷命运在史料中保存较多，城市内外仍留下的人群则难以逐一追踪。有人随军北去，有人逃往南方或附近州县，有人仍在新占领下的土地上耕作、贸易或等待亲属消息。不能因为北宋皇室被迁就把整个华北写成空地，也不能把“遗民”当作一个无差别的身份。金朝随后如何接收地方、居民如何面对新的征发和法律，属于下一章须以金代材料回答的问题。

靖康的制度后果也不是一夜清零。东京的城墙、水道、市场、仓储和文书传统仍在物质上或制度上留下痕迹，只是服务对象和权力关系改变。南方逃亡与重新建国的人们会带走北宋的官制、财政经验和记忆，同时必须在不同的地理和战争条件下改造它们。把靖康只当作北宋“失败”的道德结论，反而会遮蔽这个更重要的问题：一套曾能组织城市、财政和边防的制度，在外部征服与内部资源危机中如何分裂、迁移和被重新使用。

\subsection{不要让结局吞没前一百六十七年}

北宋从960年至1127年的历史并不是通往靖康的一条单线。陈桥后的军权重组、南方政权的接收、宋辽盟约、西夏边防、益州交子、庆历和熙宁的改革、开封的市场和徽宗朝的运输，都在各自时期解决过实际问题，也留下了新的成本。开封集中的供给系统让中枢能够支配大规模军队，也使北方道路受阻时城市格外脆弱；盟约让宋辽长期避免全面战争，也没有消除燕云问题；新法试图调整财政和基层接口，又使地方执行与派别政治更加尖锐。

这些关系不构成单一的“亡国原因”。金朝的兴起、宋金联盟的转换、河北的战争和东京的围城，是多个政权、道路和资源在特定时刻相互作用的结果。读者应把靖康视为一次关系重组的节点：北方进入金朝治理，南方出现新的宋廷，辽、西夏和东北诸区域也在不同的力量之间重排。北宋的结束因此是下一段多政权历史的开端，而不是整部故事的终点。

\section{在制度缝隙中生活：家庭、迁徙与地方社会}

朝廷材料常把北宋社会写成户口、税额、军额、学校名额或案件条目。这些分类是治理所需，却不等于生活本身。一个“户”内可能有不同年龄、性别、职业和迁徙经历的人；一条田宅文书能说明一次买卖或继承，不能代表一地全部土地关系；一方墓志保存的往往是家族希望留下的形象，不能把精英的自述当作普遍社会。北宋地方社会的困难，不是材料太少便可以用想象补足，而是材料类型彼此不相称，须让每一种材料只承担它能承担的判断。

\subsection{户籍、保甲和谁被写进名册}

户籍使国家能够分配税役、追查逃亡、安排赈济和征发兵源。名册中的户等、丁口和籍属看起来清楚，实际却不断受到出生、死亡、婚姻、分家、迁居和隐漏的影响。州县官与里正需要把这些变化写进可计算的格式，家庭则可能因税役、亲属照应或市场机会而调整居住和登记。没有一份覆盖全国、逐户连续的北宋户籍留存，故我们不能精确重建所有人的流动；但制度反复强调检括、催科与逃户，本身说明登记并不是自然完成的事实。

保甲法在熙宁后试图将若干家庭编为互保、训练和治安单位。条文能够告诉我们国家希望谁互相监督、何时操练，不能证明每个保甲都按设计运作。富户、乡里头面人物、胥吏与普通户主在地方的资源不同，执行保甲可能成为共同防卫，也可能成为新增差役和控制。用“基层组织化”一词概括，会把这些差异抹掉。更好的读法是把保甲看作国家试图缩短同乡村距离的一次尝试，其效力仍受人手、地方权力与财政条件制约。

\subsection{婚姻、继承和妇女能见度}

北宋妇女在官方本纪中很少作为独立行动者出现，却并不意味着她们不参与财产、婚姻、宗教和家庭决策。墓志会记载某些精英妇女的籍贯、婚配、子女与德行，契约和法律材料偶尔使嫁资、继承、收养和诉讼变得可见；这些证据的保存都具有选择性。墓志偏好家族的规范语言，法律案件偏向发生争执的家庭，不能从一例推及所有地区或阶层。

婚姻也是资源重新组合的方式。家庭通过嫁资、田宅、劳动和亲属网络应对风险，国家则从户籍和法律角度定义某些关系。两种视角可能重叠，也可能冲突。若只把婚姻写成礼教的实施，便看不见它如何关系到迁居、债务和照料；若只把一份契约看作自由选择，也会忽略家庭、经济和法律约束。材料允许我们讨论具体文书中的人如何主张权利，却不能替沉默者补写想法。

\subsection{流动者：道路连接的不只是货物}

汴河、长江和海路常以漕运、贡物和军粮的名义进入史书，但船上和道路上还有船户、脚夫、商旅、僧人、应试者、役夫和逃难者。国家依赖他们维持运输，又通过关津、税卡、户籍和治安试图分类他们。道路越繁忙，地方越难只凭原籍管理人群；一座桥、一个渡口、一处旅店或寺院，可能成为信息、交易和暂住相遇的场所。笔记和地理书能提供局部景象，不能转换为所有旅行者的统计。

军事危机尤其加速流动。北方战事、河患、赋役和市场波动都可能使人离开原地；南迁或投靠亲属的记录并不自动说明所有人都作出同样选择。对官府而言，流动者可能是待安置、待征发或待稽查的人；对家庭而言，他们可能是带来收入、风险或消息的亲属。把流民写成被动的灾难符号，会忽略他们在有限条件下的策略；把迁徙写成自由的市场选择，又会忽略战争和制度带来的强迫。

\subsection{地方司法：一纸判决能看见什么}

宋代法律与会要材料保存了大量关于刑名、诉讼和行政程序的规范。它们最适合说明国家期待地方官如何受理、审问、上报和判决，不足以证明所有案件都如此办理。诉讼进入档案，已意味着争端被翻译为国家认可的语言；未能告状、选择私下调解或根本不被记录的人，则较难被看见。司法材料因此不是地方社会的全景，却能显示田宅、债务、婚姻、伤害和役使如何在特定时刻成为可争辩的权利。

地方官的裁断也受时间和资源限制。案件要有文书、证人、差役与审限，州县还同时承担税收、治安和赈济。清官故事以个人道德解释裁断，容易把制度条件删去；反过来，制度条文也不能消除官员、胥吏和地方势力的裁量。北宋的地方秩序正是在这些并不完美的程序中维持、拖延或崩解。它不能用“法制发达”或“吏治败坏”一个词收束。

\section{海港与内河：北宋的市场面向哪里}

北宋的经济空间不止是东京向外辐射。四川、两浙、福建、岭南、河北和西北各有不同的物产、道路和交易对象；国家的税收、盐茶专卖和转运使这些区域相连，也不能使它们变成同一种市场。钱币、瓷器、港口遗址和船货能显示生产与流通的轨迹，官府文书能显示征税与禁令，地理笔记则记录旅行者的观察。三者回答的问题不同，不能相互替代。

\subsection{盐、茶与国家的接口}

盐和茶因日常需求广、产地集中或运输可控，成为北宋财政关心的商品。专卖制度的条目显示国家怎样设想征收、发卖和管理，却不能说明私贩是否因此消失，也不能将官方收入直接换算成生产者或消费者的负担。盐场、茶山、河道和市场之间有许多中间环节：采制者、运夫、商人、仓吏与地方官都可能影响一批货的价格和去向。对边防地区而言，盐茶还可能同军需、互市和外交相连。

制度的反复调整说明国家在追求收入、供应和控制之间寻找平衡。限制太严，可能推动绕道或私下交易；放得太宽，又可能损失财政或影响边地供应。现存材料往往更容易记录官府认为的问题，而不是所有人的交易技巧。我们可以看到国家如何试图让商品成为可征管的资源，却不能以此断言市场只按国家意志运行。

\subsection{福建、明州与海上距离}

福建和两浙的港口把北宋接到东南海域。海船、瓷器、香料和外来商人的记录使后人容易直接说“海贸繁荣”，但一次船货、一个港口遗址或一条市舶规定只能照亮有限的航线和时段。海上贸易依赖季风、造船、港湾、商人信用和地方治安，朝廷的市舶管理只是其中一环。对开封而言，海路是可增加税收和物资的远方接口；对沿海居民而言，它是渔业、运输、手工业和跨地风险交织的生活空间。

港口的多样性也不等于所有人拥有相同机会。能够留下名字的常是官员、富商、寺院捐施者或外来使者，船工和小贩较少有直接文字。考古发现能补充器物与生产，却不能从一处码头推出整个海域的政治控制。将港口写入北宋史，并不是要把它变成“开放”的装饰，而是要说明财政、城市和地方社会如何在陆路边防之外拥有另一套受天气与市场制约的联系。

\subsection{物件的旅程与证据的限度}

一件瓷器从窑场到城市或海外，可能经过窑工、装卸者、商人、船户和消费者；遗址出土却只留下它在某一处停留的最后证据。钱币窖藏同样有封藏时间、铸造时间与流通时间的差别，不能将发现地点直接等同于生产地或商业中心。物质遗存的价值不在于替文字提供一个更“客观”的结论，而在于迫使我们询问文本没有记录的生产、运输和使用环节。

北宋国家以税、专卖和文书穿过这些环节，始终无法完全取代它们。市场和城市之所以能支撑财政，是因为无数地方性的劳动和信任关系在运作；国家试图将之分类和征收，又不断受到自然条件、信息不全和地方协商的限制。理解这一点，才不会把北宋经济写成一台由东京单独操纵的机器。
